\scanpage{317}
\mspage{16}
(c’est ce qui se vérifie directement dans certains cas remarquables, p.ex. celui où $L=\varnothing$ et où on trouve la\unsure{\inserted{\unsure{\scriptsize q}}} catégorie pleine de $\mathrm{Ab}_*$ engendrée par \inserted{les} $\unsure{G_o[n]|C_o[n]}$ grâce aux seules sommes finies.

c) On a oublié de noter que \deleted{dans} dans le cas particulier $B=\mathrm{Ab}$, on trouve les $=$ types d’objets (savoir ceux de $\mathrm{Ab}_*$ ou de $\mathrm{Ab}_\bullet$ au choix) pour exprimer les différents \uline{inv.} contravariants \inserted{$F$} possibles à valeurs dans $\mathrm{Ab}$ (commutant aux $\varprojlim$) et les \uline{op.} \inserted{$\Omega$} qu’on peut faire sur des inv. contravar. \deleted{quelconques} à val. dans une cat. \deleted{\unsure{munie}} \inserted{additive} à $\varprojlim$, On a déjà dit comment
\[
\diagram{\begin{array}{ccl}
F&\text{correspond à}&K_*\in\mathrm{Ab}_*\quad\text{ou}\quad\unsure{k_\bullet|K_\bullet}\in\mathrm{Ab}_\bullet\\
\Omega&\unsure{\longleftarrow}&L_*\in\mathrm{Ab}_*\quad\text{ou}\quad L_\bullet\in\mathrm{Ab}_\bullet
\end{array}}
\]
Comment alors exprimer directement la correspondance entre $F$ et $\Omega$, si $K=L$ ? On trouve si $F_o\in((\mathrm{Ab}_*)^o)_*=B_{o*}$ \inserted{$\simeq({\mathrm{Ab}_*}^{*})^o$} est le \deleted{\unsure{covariant}} \inserted{l’invariant} universel comme dessus, à valeurs dans $B_o=(\mathrm{Ab}_*)^o$, alors pour tt $X_*\in\mathrm{Ss}=\widehat{\boldsymbol\Delta}$
\[
\boxed{F(X_*)\simeq\Omega_{\unsure{*}}^*(C_*(X_*)^o)\simeq\Omega_{\unsure{\bullet}}^\bullet(C_\bullet(X_*)^o)}
\]
\margin{318}

où l’exposant $o$ à $C_*(X_*)$ ou $C_\bullet(X_*)$ signifie les ms, mais considérés comme objets non de $\mathrm{Ab}_*$ \inserted{resp $\mathrm{Ab}_\bullet$}, mais de \deleted{\unsure{(\mathrm{Ab})^o}} ${B'}^*$ \inserted{resp ${B'}^\bullet$} où $B'=\mathrm{Ab}^o$, de sorte que $\Omega^*$ \deleted{resp. $\Omega^\bullet$} est défini sur \deleted{$\mathrm{Ab}_*\ \text{resp}\ A$} ${B'}^*$ resp ${B'}^\bullet$ ($\Omega_B^*:{B'}^*\longrightarrow\unsure{B'|B^\vee},\ \Omega_B^\bullet:{B'}^\bullet\longrightarrow B'$) a un sens sur $C_*(X_*)^o$ \inserted{resp. $C_\bullet(X_*)$}, et c’est cet élément de $B=(\mathrm{Ab})^o$ qui est $F(X_*)$ (à isom. can. près).

Tous ces développements sont encore valables pour une catégorie $\Delta$ (ess. petite) quelconque --- et aussi dans le contexte non additif --- \deleted{\unsure{dans ce qui est}} \unsure{sauf les} $*$ et \unsure{ceux en} $\bullet$, qui, eux, nécessitent des catégories de valeurs additives et la théorie bien spéciale de Dold-Puppe.

En fait (revenant au paquet $L$ de cat. d’indices envisagé dans b)) on s’intéressera surtout aux inv. contravariants en $X_*$ provenant d’objets $K_*\in\mathrm{Ab}_*$ qui sont dans $U_L$ --- i.e. qui peuvent se déduire des complexes $C_*^{[n]}\in(\mathrm{Ab})_*$ (et des opérations $\tfrac12$ simpliciales entre lesdites), à l’aide des \unsure{opérations} $\Omega_*$ \deleted{\unsure{universelles}\ \unsure{de}\ \unsure{\mathrm{Ab}_*}} \inserted{\unsure{provenant par là}, soit des $\unsure{C_*^n}\in(\mathrm{Ab})_*$ et des op. de $\unsure{\tfrac12}$ \unsure{simplicialité}, soit des \unsure{telles opérations}} \unsure{ce qui revient au $\subset$ de dire que l’on s’intéressera aux foncteurs} \deleted{$\unsure{F(X_*)}$}

\margin{On pourrait comme d’habitude, introduire un $h$, remplaçant partout $\mathrm{Ab}$ par $\mathrm{Ab}_h$ ---}
\margin{(NB) attention, dans cette formule \unsure{apparaît la confusion en attendant} de la situation --- il s’agit bien des $F$ \uline{contravariants} à val. dans $\mathrm{Ab}$ lui-même (et non dans $(\mathrm{Ab})^o$ \unsure{---}).}
